The results say that the separate treatment of the UF changes the Annex 1 measures in two different ways, and that the two ways have different causes. The income measure differs from its single-currency counterpart through two additive terms that need no market input: the charge, which is a fixed fraction of the position, and the offset that is only half of what a netting would give. Both are rules of the text and not estimates, and the proportion of about one to two in which the pair rule and the charge carry the gap is stable across books (Table~\ref{tab:shapley}). The value measure differs through two inputs that depend on the market and on the supervisor's calibration, the real curve and the table of sizes, whose effects nearly cancel on the book studied. The sign of each is as Proposition~\ref{prop:curve} says, a large positive effect of the real curve on the sensitivity of the UF flows against a negative one of the smaller sizes, but the near cancellation is a property of the book and of the two curves of the valuation date, not of the treatment: the weight $w_k$ crosses one between \tStarLo{} and \tStarHi{} years, and a book whose UF flows lay beyond that horizon would have the single-currency treatment below Annex 1 for the value measure by more than the book studied, while a book of short flows would have it above.

The failures are informative. H2 failed because it expected the aggregation rule, the visible feature of the text, to carry the value measure; the data say the inputs carry it. H3 failed in one clause because the charge is a constant that does not scale with the size of the offset, so that for a small net UF position the Annex 1 income measure can fall below the Basel positive-part aggregation, a condition that Equation~\eqref{eq:charge-condition} states exactly and that makes the Basel aggregation, in books with a large offset and a small position, the more conservative of the two. H5 failed because a floor, which is natural to read as a limit on the shock, interacts with a pair rule that offsets gains: reducing the gain of one leg reduces what is offset, and the measure rises. A supervisor or a bank that tests the floor by asking whether the measure falls when it is imposed would be misled in the books that Table~\ref{tab:floors} shows.

The practical reading depends on the audience. For a bank, the charge is a lower bound on the income measure that no hedge of the rate risk removes, and it is 2.5 percent of the position, so that a net UF position of 30 percent of level 1 capital uses 0.75 percent of the 5 percent line (\repChargeOverTier{} percent in the short UF variant). The test of the income measure cannot be breached by the charge alone short of a position of \floorBreakNu{} times capital, so for a bank with a UF position of ordinary size the measure that classifies it is the value measure, which the near cancellation of the two inputs leaves close to what a single-currency calculation gives (the single-currency value measure is $\mainSingleOverAnnexEveLng$ of Annex 1 in the long UF variant). For a validator of such a bank's model, the decomposition gives a way to test an implementation: the dummy and the charge of Proposition~\ref{prop:shapley} are exact, the sandwich of Proposition~\ref{prop:pair} is an inequality that any implementation of the pair rule must satisfy, and a defect in either shows at once.

The result of the sweep, that the two treatments classify more than half of the books differently only at positions of \sweepMajLo{} to \sweepMajHi{} percent of level 1 capital, the top of the range tested, and single books from \sweepFirstDiff{} percent, has a second reading. A treatment that differs by a factor of about five in an income measure and by \mainEveShortLo{} to \mainEveShortHi{} percent in a value measure, and yet changes the classification of no book at the positions of 30 percent and below, is a difference that is invisible to the thresholds and visible in the limits. The board limits of the text are set as percentages of the margin and of level 1 capital, and a limit on the income measure at a fixed share of the margin would be reached at very different positions under the two treatments. The classification is what the supervisor uses, so the factor measures a difference in the measure and not in the supervisory outcome at the positions studied; the paper therefore reports what each component adds to the measures and goes no further.
